\section{Равновесные составы при 1000 К}
\label{sec:bulk-results}

При нулевом внешнем давлении, свободных объёме и отношении осей
$c/a$ прямой расчёт на 1400 атомах дал
\[
 x_{\alpha}=14{,}982\ \text{ат.\% Al}, \qquad
 x_{\alpha_2}=22{,}445\ \text{ат.\% Al}.
\]
Для этой оценки выбран более широкий из двух проверенных пределов
смещений, $W<0{,}40$. Условные 95\%-е интервалы по блокам траекторий
равны $14{,}917$--$15{,}063$ и $22{,}327$--$22{,}553$ ат.\% Al.
Ширина интервала сосуществования составляет $7{,}463$ процентного
пункта (п.п.), с 95\%-м интервалом $7{,}323$--$7{,}580$ п.п.
Здесь и далее составы даны в атомных процентах Al, а изменения
составов --- в процентных пунктах.

\subsection{Размер ячейки и статистическое разрешение}

Для каждого из четырёх размеров изучены две области допустимых
смещений. Вся серия охватывает десять физических ансамблей при
различных $N$, пределе $W$ и условиях на $c/a$. В них накоплено
$255\,200$ производственных проходов в 268 зонтичных траекториях.
Ещё 16 независимых траекторий дали $11\,200$ проходов без зонтичного
ограничения. Один проход включает $N$ попыток химической замены
и $N$ попыток локального смещения; изменение объёма предлагается
каждый проход. Первоначальная релаксация исключена из всех средних.

\begin{table}[htbp]
\centering\small
\caption{Составы фаз (ат.\% Al) при свободном $c/a$. Для более
широкой области приведён также барьер $\Delta F_Q$ между максимумами
распределения $Q$, измеренный в поле равных высот.}
\label{tab:bulk-size}
\begin{tabular}{@{}r c rr rr r@{}}
\toprule
 & & \multicolumn{2}{c}{$W<0{,}35$} & \multicolumn{2}{c}{$W<0{,}40$}
 & $W<0{,}40$\\
\cmidrule(lr){3-4}\cmidrule(lr){5-6}
$N$ & Ячейка & $\alpha$ & $\alpha_2$ & $\alpha$ & $\alpha_2$
 & $\Delta F_Q/(k_{\mathrm B}T)$\\
\midrule
  96 & $4\times4\times3$   & $14{,}93$ & $22{,}83$ & $15{,}48$ & $22{,}86$ & $1{,}18$\\
 288 & $6\times6\times4$   & $15{,}04$ & $22{,}47$ & $15{,}19$ & $22{,}42$ & $2{,}17$\\
 768 & $8\times8\times6$   & $15{,}02$ & $22{,}52$ & $14{,}93$ & $22{,}40$ & $4{,}23$\\
1400 & $10\times10\times7$ & $14{,}86$ & $22{,}36$ & $14{,}98$ & $22{,}44$ & $5{,}66$\\
\bottomrule
\end{tabular}
\end{table}

При $W<0{,}40$ увеличение $N$ с 768 до 1400 сдвинуло обе границы
на $+0{,}05$ п.п. Условные 95\%-е интервалы этих разностей равны
$[-0{,}05;\,+0{,}16]$ для $\alpha$ и $[-0{,}11;\,+0{,}21]$ п.п.
для $\alpha_2$. При $W<0{,}35$ на том же шаге получены
$-0{,}166$ и $-0{,}159$ п.п. Для $\alpha$ при более тесном
ограничении интервал $[-0{,}298;\,-0{,}034]$ п.п. указывает на
ещё различимый размерный сдвиг. Поэтому в качестве результата
используется непосредственно измеренная точка при $N=1400$.
По этим размерам и независимым стартам практически достижимое
разрешение сравнения составляет около $0{,}2$ п.п. Бутстрэп
выполнялся по блокам внутри каждой траектории с повторной оценкой
поля и составов; при $N=768$ и 1400 использовано по 192 реализации.
Такой интервал отражает неопределённость выборки внутри заданного
ансамбля. Влияние размера, формы и ограничения смещений оценивается
отдельными расчётами.

При $N=1400$ и $W<0{,}40$ поле сосуществования равно
$\Delta\mu_{\mathrm{cfg}}=3{,}29368$ эВ на замену Ti$\to$Al
в конфигурационной нормировке. Максимумы распределения $Q$ при
этом поле находятся около $0{,}063$ и $0{,}878$. В фазе $\alpha_2$
на наиболее богатой алюминием подрешётке Al занимает в среднем
$87{,}5\%$ узлов; на остальных трёх, упорядоченных по убыванию
заполнения, --- $1{,}19\%$, $0{,}72\%$ и $0{,}37\%$.
Это характерный признак порядка D0$_{19}$ при отклонении от
стехиометрии Ti$_3$Al. Объём на атом и $c/a$ составили
$17{,}457$~\AA$^3$ и $1{,}5915$ для $\alpha$, $17{,}241$~\AA$^3$
и $1{,}6064$ для $\alpha_2$. В пределах выбранных диагностических
полос у границ допустимого объёма и $c/a$ статистический вес обеих
фаз оказался нулевым.

\subsection{Чувствительность к ограничению смещений}

При $N=1400$ расширение области с $W<0{,}35$ до $W<0{,}40$
повысило поле сосуществования на $4{,}59$ мэВ на замену
Ti$\to$Al. Состав $\alpha$ изменился на $+0{,}123$ п.п.
с 95\%-м интервалом от $+0{,}003$ до $+0{,}251$ п.п.
Для $\alpha_2$ сдвиг составил $+0{,}083$ п.п. с интервалом
от $-0{,}079$ до $+0{,}251$ п.п.
Более широкая область принимает дополнительные состояния с
большими тепловыми смещениями. При $N=288$ в ансамбле $W<0{,}40$
вероятность выполнить более строгое условие $W<0{,}35$ составила
$0{,}158$ в $\alpha$ и $0{,}468$ в $\alpha_2$. Поэтому ужесточение
ограничения термодинамически благоприятствует упорядоченной фазе.
Повторное взвешивание отобранных конфигураций дало
$15{,}05$ и $22{,}34$ ат.\% Al. В самостоятельном расчёте
при $W<0{,}35$ получены $15{,}04$ и $22{,}47$ ат.\%.

Условие на $W$ должно выполняться одновременно на всех атомах.
На 1400 атомах вероятность пройти эту же фильтрацию в фазе $\alpha$
упала до $1{,}23\cdot10^{-5}$. Эффективное число сохранившихся
кадров по весам Киша составило всего $3{,}6$ ещё до поправки на
автокорреляцию. Для наибольшей ячейки вывод о чувствительности
основан на самостоятельных сериях при каждой величине $W$.
Обе серии характеризуют конденсированный ГПУ-бассейн данного MTP.

\subsection{Проверки положения границ}

Два независимых набора зонтичных траекторий на 1400 атомах при
$W<0{,}40$ дали для $\alpha$ значения $14{,}89$ и $15{,}08$
ат.\%, а для $\alpha_2$ --- $22{,}37$ и $22{,}53$ ат.\%.
При переносе разделяющей поверхности от $Q=0{,}34$ до $0{,}54$
каждая граница меняется максимум на $0{,}02$ п.п. Если в той же
ячейке потребовать равенства \emph{суммарных} весов бассейнов,
получаются $14{,}90$ и $22{,}34$ ат.\%; основной расчёт учитывает
вырождение четырёх вариантов порядка в $\alpha_2$.

Две проверки использовали распределение состава без деления по $Q$.
Пересечение зависимостей $\langle x\rangle(\Delta\mu)$ для
$N=768$ и 1400 при $W<0{,}40$ дало
$3{,}29397$ эВ с интервалом $3{,}29235$--$3{,}29640$ эВ.
Вторая проверка использует четвёртую кумулянту состава
$\kappa_4=\langle(x-\langle x\rangle)^4\rangle-
3\,\mathrm{Var}(x)^2$ в максимуме дисперсии. Для равновесной смеси
двух одинаково взвешенных узких пиков расстояние между ними
оценивается как $(-8\kappa_4)^{1/4}$. Эта оценка дала
$7{,}48$ и $7{,}46$ п.п. при $N=768$ и 1400 соответственно.
Она подтверждает ширину по разности фазовых средних, $7{,}46$ п.п.,
с помощью другого наблюдаемого на тех же выборках.

В независимом тесте для каждой величины $W$ запущено по четыре
случайных и четыре упорядоченных состояния на 1400 атомах.
По химическому порядку они эволюционировали свободно. Для каждой
фазы состав оценивался отдельно с поправкой на разницу между
заданным полем и полем сосуществования из зонтичных окон.
При $W<0{,}40$ получены $14{,}987$ и $22{,}489$ ат.\% Al,
с 95\%-ми интервалами по репликам и полю $14{,}88$--$15{,}11$
и $22{,}29$--$22{,}64$ ат.\%. При $W<0{,}35$ получены
$14{,}914$ и $22{,}486$ ат.\%; две первоначально упорядоченные
реплики перешли в $\alpha$ во время термализации. В производственной
части переходов между ядрами бассейнов не обнаружено.
Эти траектории проверяют \emph{составы фаз при найденном поле};
само поле определено из относительных статистических весов зонтичных
окон.

\subsection{Сопоставление с C46 при согласованной геометрии}
\label{sec:bulk-c46}

Для C46 повторно восстановлена общая касательная по 13 авторским
окнам при 1000 К, $p=0$ и $N=5120$. Составы совпали с исходным
результатом до точности записи. Геометрический контроль прямого
расчёта сохранял отношение осей C46, $c/a=4{,}6046/2{,}9446
=1{,}56374$, при свободном объёме. Для этого контроля поле
сосуществования уточнялось независимо, включая две серии окон
непосредственно вблизи $3{,}322$ эВ.

\begin{table}[htbp]
\centering\small
\caption{Геометрия ячейки и составы. Для MTP здесь $N=1400$,
$W<0{,}40$. В квадратных скобках приведены условные 95\%-е
статистические интервалы. Ширина $\Delta x=x_{\alpha_2}-x_{\alpha}$
дана в п.п.}
\label{tab:bulk-c46}
\begin{tabular}{@{}l c c r@{}}
\toprule
Метод и форма & $\alpha$, ат.\% Al & $\alpha_2$, ат.\% Al & $\Delta x$\\
\midrule
MTP, свободное $c/a$ & $14{,}98\ [14{,}92;\,15{,}06]$ & $22{,}44\ [22{,}33;\,22{,}55]$ & $7{,}46$\\
MTP, $c/a$ как в C46 & $16{,}08\ [16{,}03;\,16{,}14]$ & $22{,}16\ [22{,}06;\,22{,}24]$ & $6{,}08$\\
C46, заданное $c/a$ & $14{,}96\ [14{,}93;\,14{,}99]$ & $21{,}96\ [21{,}91;\,22{,}01]$ & $7{,}00$\\
\bottomrule
\end{tabular}
\end{table}

В фиксированной геометрии увеличение ячейки с 768 до 1400 атомов
сдвинуло оценки на $-0{,}036$ и $-0{,}065$ п.п. Использование
только двух серий вблизи поля сосуществования дало
$16{,}091$ и $22{,}169$ ат.\% при $N=1400$. Это согласуется
с объединёнными оценками $16{,}083$ и $22{,}159$ ат.\%.
Поле в конфигурационной нормировке равно $3{,}32224$ эВ.

Сравнение с согласованным $c/a$ даёт совокупный сдвиг от C46
к прямому MTP $+1{,}12$ п.п. для $\alpha$ и $+0{,}20$ п.п.
для $\alpha_2$. При последующей релаксации отношения осей в прямом
расчёте граница $\alpha$ сдвигается на $-1{,}10$ п.п.,
а граница $\alpha_2$ --- на $+0{,}29$ п.п. Их 95\%-е интервалы
составляют от $-1{,}19$ до $-1{,}00$ и от $+0{,}14$ до
$+0{,}43$ п.п. соответственно.
В итоге относительно C46 получаем $+0{,}02$ и $+0{,}49$ п.п.
Близость двух конечных оценок для $\alpha$ возникает вследствие
компенсации сдвигов противоположного знака. Ширина области
сосуществования при освобождении осей возрастает на $1{,}39$ п.п.
с интервалом $1{,}22$--$1{,}54$ п.п.

У обоих расчётов один исходный MTP. При совпадающем $c/a$ остаются
различия в представлении химической энергии и колебаний
кластерным разложением C46, в квазигармоническом приближении и
в статистике ядер (квантовая в C46, классическая в прямом
расчёте). Измеренная разность характеризует их совокупный вклад.

\paragraph{Прослеживаемость.}
На UMT под учётной записью \texttt{u3129} завершены 26 производственных
серий и независимых контролей. Исходные идентификаторы Slurm,
SHA-256 архивов и число проходов по сериям сохранены в
\texttt{run\_provenance.json}. Числа для таблиц получены из
\texttt{summary\_statistics.json} и \texttt{c46\_reference.json},
проверки по составу и прямым траекториям ---
из \texttt{phase\_checks.json} и \texttt{endpoint\_controls.json}.
Вычислитель энергии отдельно сопоставлен с MLIP3 на 123 конфигурациях
из 8, 128 и 432 атомов, включая смены типа, смещения и изменения
объёма. Максимальное абсолютное расхождение составило
$1{,}36\cdot10^{-12}$ эВ на всю ячейку.
